\documentclass[11pt,a4paper]{article}

% ── Packages ──
\usepackage[utf8]{inputenc}
\usepackage[T1]{fontenc}
\usepackage{amsmath,amssymb}
\usepackage{graphicx}
\usepackage[margin=1in]{geometry}
\usepackage{natbib}
\usepackage{hyperref}
\usepackage{booktabs}
\usepackage{xcolor}
\usepackage{siunitx}
\usepackage{caption}
\usepackage{subcaption}
\usepackage{lineno}

\linenumbers

% ── Macros ──
\newcommand{\vmr}{\mathrm{VMR}}
\newcommand{\um}{\,\si{\micro\meter}}
\newcommand{\chisq}{\chi^2}

% ── Title ──
\title{Multiscale Spatial Fingerprints Reveal Cell Type Organization\\in the Mouse Brain}

\author{
  Pretham Voruganti\textsuperscript{1}
  \\[0.5em]
  \small\textsuperscript{1}Independent Researcher
}

\date{}

\begin{document}

\maketitle

% ══════════════════════════════════════════════════════════════
\begin{abstract}
Spatial transcriptomics reveals the positions of thousands of cell types across
tissue, but current analysis tools either operate at a single spatial scale or
lack interpretable statistical summaries. Here we introduce a multiscale spatial
fingerprinting framework that quantifies how cell types are organized across
spatial scales from \SI{40}{\micro\meter} to \SI{1000}{\micro\meter} in the
mouse brain. Using Allen Brain Cell Atlas MERFISH data
(\num{1374804} cells across isocortex, hippocampus, and thalamus), we compute
variance-to-mean ratio (VMR) fingerprints for each cell type at 12 logarithmically
spaced scales, tested against label-shuffle null models via covariance-aware
Hotelling's $T^2$ statistics.

We show three principal results.
First, cell types exhibit robust, grid-placement-independent spatial fingerprints
whose effect sizes vary over two orders of magnitude (peak VMR ratios
\numrange{1.2}{36.4}$\times$ null).
Second, subclass-level fingerprints in isocortex recover known laminar
architecture, with excitatory subtypes showing layer-specific characteristic
clustering scales, while thalamic subtypes show distinct nuclear
compartmentalization signatures.
Third, cross-type co-occurrence is scale-dependent: isocortical layer types
predominantly co-localize (96--98\% positive at all scales), whereas thalamic
nuclear types show mixed, weak associations (49--57\% positive), consistent
with the distinction between layered and nuclear brain architectures.

The same VMR/chi-squared framework previously applied to protein nanoclusters
(SMLM, nm scale) and synaptic clustering on dendrites (\si{\micro\meter} scale)
here detects meaningful structure at the tissue scale (\si{\milli\meter}),
supporting a universal spatial fingerprinting approach across biological scales
and geometries.
\end{abstract}

\clearpage
\tableofcontents
\clearpage

% ══════════════════════════════════════════════════════════════
\section{Introduction}
\label{sec:intro}

% ── The problem ──
Spatial transcriptomics technologies---including MERFISH, seqFISH+, Slide-seq,
and Visium---now routinely resolve the positions and transcriptomic identities of
millions of cells across tissue sections \citep{Zhang2023}. A central question
emerging from these datasets is: \emph{how are cell types spatially organized,
and at what scales does that organization manifest?}

% ── Existing tools and their limitations ──
Current spatial analysis tools address parts of this question. Squidpy
\citep{Palla2022} computes neighborhood enrichment at a fixed spatial
neighborhood radius. SpatialDE \citep{Svensson2018} identifies spatially variable
genes but operates gene-by-gene, not type-by-type. BANKSY \citep{Singhal2024}
incorporates spatial context into clustering but does not provide
scale-resolved statistics. None of these tools produce continuous
\emph{multiscale spatial fingerprints}---curves that summarize how a cell type's
clustering varies continuously from fine ($\sim$\SI{40}{\micro\meter}) to coarse
($\sim$\SI{1000}{\micro\meter}) resolution.

% ── Our approach ──
Here we apply a multiscale variance-to-mean ratio (VMR) framework, previously
developed for protein nanocluster detection in super-resolution microscopy and
synapse clustering on dendritic trees, to the tissue scale. The key idea is
simple: at each spatial scale, tile the tissue into non-overlapping grid cells,
count instances of a cell type per cell, and compute the VMR of those counts.
For a completely random (Poisson) distribution, VMR~$= 1$; clustering produces
VMR~$> 1$. By sweeping across scales, we obtain a \emph{spatial fingerprint}
that encodes where and how strongly a cell type clusters.

% ── Null model ──
Significance is assessed via label-shuffle null models---the strongest available
control, which preserves all spatial structure and only permutes type-to-position
assignments. The full VMR curve is tested jointly using a covariance-aware
Hotelling's $T^2$ statistic with Ledoit--Wolf shrinkage regularization.

% ── Scope ──
We apply this framework to three brain regions from the Allen Brain Cell Atlas
MERFISH dataset: isocortex (\num{936281} cells), hippocampal formation
(\num{304736} cells), and thalamus (\num{133787} cells). We analyze per-type
fingerprints, cross-type co-occurrence, regional comparisons, and a focused
comparison of layered (isocortex) versus nuclear (thalamus) organizational
motifs.


% ══════════════════════════════════════════════════════════════
\section{Methods}
\label{sec:methods}

\subsection{Dataset}
\label{sec:dataset}

We used the Allen Brain Cell Atlas MERFISH dataset (MERFISH-C57BL6J-638850)
\citep{Zhang2023}, which provides spatial coordinates and transcriptomic cell
type annotations for \num{3938808} cells across the adult mouse brain. Cells
are annotated at three hierarchical levels: 34 classes, 338 subclasses, and
1,201 supertypes.

Brain region assignments were obtained by merging cell coordinates with the
Allen Common Coordinate Framework (CCF v3) parcellation. We analyzed three
regions: isocortex (CCF division ``Isocortex''), hippocampal formation (``HPF''),
and thalamus (``TH'').

\subsection{Multiscale VMR Fingerprints}
\label{sec:vmr}

For a given cell type $t$ with positions $\{(x_i, y_i)\}_{i=1}^{n_t}$:

\begin{enumerate}
  \item Construct a 2D density grid ($512 \times 512$ pixels) over the region
    of interest via histogram binning.
  \item For each spatial scale $s$ in a set of 12 logarithmically spaced
    scales from \SI{40}{\micro\meter} to \SI{1000}{\micro\meter}:
    \begin{enumerate}
      \item Partition the grid into non-overlapping cells of side length $s$.
      \item Count the number of type-$t$ cells per grid cell.
      \item Compute $\vmr(s) = \mathrm{Var}(c) / \mathrm{Mean}(c)$.
    \end{enumerate}
  \item The resulting curve $\vmr(s)$ is the \emph{spatial fingerprint} of type~$t$.
\end{enumerate}

Scale selection uses the \texttt{ScaleRange.for\_tissue()} algorithm, which
generates candidate scales in log-space, converts to grid voxel sizes, deduplicates
to ensure each scale produces a distinct binning resolution, and enforces a
minimum resolvable scale of $2 \times (\text{ROI size} / \text{grid size})$.

\subsection{Null Models and Statistical Testing}
\label{sec:null}

\paragraph{Label-shuffle null.}
For each of 50 realizations, cell type labels are randomly permuted while
positions remain fixed. This preserves all spatial structure (density gradients,
tissue geometry) and tests specifically whether type-to-position assignments are
non-random.

\paragraph{Covariance-aware chi-squared test.}
The observed VMR curve $\mathbf{v}_\text{real}$ is compared to the null
distribution $\{\mathbf{v}_k\}_{k=1}^{K}$ using Hotelling's $T^2$:
\begin{equation}
  T^2 = \frac{K}{K+1} (\mathbf{v}_\text{real} - \bar{\mathbf{v}})^\top
  \hat{\Sigma}^{-1}_\text{reg}
  (\mathbf{v}_\text{real} - \bar{\mathbf{v}})
\end{equation}
where $\hat{\Sigma}_\text{reg} = (1-\lambda)\hat{\Sigma} + \lambda \,
\mathrm{diag}(\hat{\Sigma})$ is the Ledoit--Wolf shrinkage-regularized
covariance ($\lambda = 0.02$). The $T^2$ statistic is converted to an
$F$-statistic with $(p, K-p)$ degrees of freedom.

\paragraph{Effect-size metrics.}
To avoid over-reliance on $p$-values in this high-powered setting, we report:
\begin{itemize}
  \item \textbf{Peak VMR ratio}: $\max_s \vmr_\text{real}(s) / \vmr_\text{null}(s)$
  \item \textbf{Integrated excess}: $\int [\vmr_\text{ratio}(s) - 1] \, d(\log s)$
  \item \textbf{Fingerprint width}: scale range (in decades) over which VMR ratio
    exceeds half its peak
\end{itemize}

\subsection{Cross-Type Co-occurrence}
\label{sec:cooccurrence}

For each pair of cell types $(A, B)$ at each scale $s$, we compute the Pearson
correlation of their binned counts across grid cells:
\begin{equation}
  r_{AB}(s) = \mathrm{Corr}\bigl(c_A^{(s)}, c_B^{(s)}\bigr)
\end{equation}
Positive values indicate co-localization; negative values indicate segregation.
Significance is assessed by comparing against the label-shuffle null distribution.

\subsection{Grid Jitter Robustness}
\label{sec:jitter}

To verify that results are not artifacts of arbitrary grid placement, we repeated
the full analysis 10 times with random grid origin offsets drawn from
$\mathrm{Uniform}(-2\Delta, +2\Delta)$ in each axis, where $\Delta$ is the
grid voxel size. We report the coefficient of variation (CV) of peak VMR ratio
across jitter trials.

\subsection{Software and Reproducibility}

All analyses were implemented in Python using NumPy, SciPy, Pandas, and
Matplotlib. Data access used the \texttt{abc\_atlas\_access} package. Code,
results, and figures are available at
\url{https://github.com/saivoru777-ship-it/spatial-transcriptomics-fingerprints}.


% ══════════════════════════════════════════════════════════════
\section{Results}
\label{sec:results}

% ── Result 1 ──
\subsection{Cell types exhibit robust multiscale spatial fingerprints}
\label{sec:result1}

% Per-type fingerprints across three regions
% Effect-size ranking (not significance counting)
% Grid jitter robustness (CV < 5% for all types)
% Native vs mislocated type distinction
% Boundary crop stability

\textbf{[Figure: fig\_per\_type\_fingerprints\_isocortex.pdf]}

\textbf{[Figure: fig\_top\_types\_by\_effect.pdf]}

\textbf{[Figure: fig\_grid\_jitter\_robustness.pdf]}

\textbf{[Figure: fig\_native\_vs\_mislocated.pdf]}

\textbf{[Figure: fig\_boundary\_sensitivity.pdf]}


% ── Result 2 ──
\subsection{Fingerprints recover known anatomical organization}
\label{sec:result2}

% Subclass-level isocortex: layer types with distinct scales
% Regional comparison: isocortex vs hippocampus vs thalamus
% Known biology recovery paragraph

\textbf{[Figure: fig\_isocortex\_subclass\_fingerprints.pdf]}

\textbf{[Figure: fig\_regional\_comparison.pdf]}


% ── Result 3 ──
\subsection{Cross-type spatial association is scale-dependent}
\label{sec:result3}

% Curated co-occurrence panels
% Layered vs nuclear topology comparison
% Cross-type heatmap (supplementary?)

\textbf{[Figure: fig\_curated\_cooccurrence\_isocortex.pdf]}

\textbf{[Figure: fig\_layered\_vs\_nuclear.pdf]}


% ══════════════════════════════════════════════════════════════
\section{Discussion}
\label{sec:discussion}

% ── Summary of findings ──
We have shown that multiscale VMR fingerprints provide a compact, interpretable
summary of cell type spatial organization in the mouse brain. Three properties
make this framework valuable.

% ── Robustness ──
\paragraph{Robustness.}
Fingerprints are stable to grid placement (CV~$< 5\%$ across 10 jitter trials)
and interior cropping. The extreme effect-size tail is dominated by cell types
that are rare visitors from adjacent brain regions, clustering tightly at their
primary anatomical location. When these are separated from region-native types,
the native effect sizes span a biologically interpretable range
(peak ratios 1.0--5.8$\times$ for most types, with the habenular MH-LH Glut
reaching 36$\times$ due to tight nuclear confinement).

% ── Biology recovery ──
\paragraph{Known biology recovery.}
Subclass-level fingerprints in isocortex resolve the expected laminar hierarchy:
L2/3, L4/5, L5, L6, and L6b excitatory subtypes each show distinct clustering
scales. Glial and vascular types show characteristic small-scale
($\sim$\SI{40}{\micro\meter}) organization consistent with local microenvironment
patterning. Hippocampal fingerprints show stronger mesoscale compartmentalization
than isocortex, consistent with the highly structured dentate gyrus and CA
subfield architecture.

% ── Layered vs nuclear ──
\paragraph{Layered versus nuclear organization.}
The subclass-level co-occurrence comparison between isocortex and thalamus
reveals qualitatively distinct organizational motifs. Cortical layer types are
overwhelmingly co-localized at all scales (96--98\% of pairs positive),
consistent with laminar stacking within shared cortical columns. In contrast,
thalamic nuclear types show weak, mixed correlations (49--57\% positive), with
more frequent scale-switching---consistent with a compartmental architecture
where distinct nuclei occupy non-overlapping territories. This suggests that the
framework can distinguish between fundamentally different modes of tissue
organization.

% ── Universality ──
\paragraph{Universality across biological scales.}
The same VMR/covariance-aware chi-squared framework has now been applied to three
distinct biological systems: protein nanoclusters at nm resolution (SMLM),
synaptic input clustering at \si{\micro\meter} resolution (dendritic trees), and
cell type organization at \si{\milli\meter} resolution (tissue sections). In each
case, the framework detects meaningful, scale-specific spatial structure with
appropriate null models. This convergence supports the concept of a universal
spatial fingerprinting toolkit for biology.

% ── Limitations ──
\paragraph{Limitations.}
Several limitations should be noted. First, in datasets of this size, nearly all
types are statistically significant, making $p$-values uninformative for
distinguishing biologically meaningful from trivial effects; we address this by
reporting effect sizes throughout. Second, Pearson correlation for co-occurrence
is simple but may miss more complex spatial relationships; future work could
incorporate mark correlation functions or cross-$K$ statistics. Third, our
analysis uses 2D projections of a 3D brain; extension to full 3D spatial
fingerprints would be valuable for datasets with sufficient axial resolution.

% ── Future directions ──
\paragraph{Future directions.}
A systematic subclass-level analysis across all brain regions---beyond the
isocortex and thalamus focus here---could test whether the layered/nuclear
distinction extends to other organizational motifs (e.g., the nuclear but
layered cerebellar cortex). Integration with gene expression profiles could
reveal whether spatial fingerprint similarity predicts transcriptomic similarity.


% ══════════════════════════════════════════════════════════════
\section*{Data Availability}

The Allen Brain Cell Atlas MERFISH dataset is publicly available at
\url{https://portal.brain-map.org/atlases-and-data/bkp/abc-atlas}. All code,
computed results, and figures are available at
\url{https://github.com/saivoru777-ship-it/spatial-transcriptomics-fingerprints}.


% ══════════════════════════════════════════════════════════════
\section*{Acknowledgments}

We thank the Allen Institute for Brain Science for making the MERFISH dataset
publicly available. This work used the \texttt{abc\_atlas\_access} Python package
for data retrieval.


% ══════════════════════════════════════════════════════════════
\bibliographystyle{unsrtnat}
\bibliography{references}

\end{document}
